\section{方法评价与阶段性结论}
\subsection{模型的适用性与局限}
本文框架把单点能力、调度、中继和分区组织在同一套物理与资源定义之下。其可解释性来自对象层次的区分：货箱是交付最小单位，架次是飞行任务单位，实体机与电池是分别占用的资源，通信分段是实际保障关系的记录单位。这样可以定位容量不足、能源等待或中继复制的具体来源。

模型的精确性范围也具有层次。单点集合划分在候选完整时可精确求解；固定时序区间模型可精确核算资源数；多点运输和静态中继候选的联合构造仍受启发式与候选集合限制。即使连续通信检验通过，也只证明该方案满足指定传播判据，并不证明中继点位或整体目标达到全局最优。

实际救援还受到天气、定位误差、临时禁飞、充电设备条件和动态需求影响。当前模型采用附件确定参数，30米DEM不能表征全部小尺度障碍，题设传播公式也不是完整无线环境仿真。因此，得到的方案应理解为标准场景下的计划依据，其推广需要额外的气象、通信和装备实测数据。

\subsection{已经形成的认识与后续定量问题}
从原始附件可以确认，运输需求同时包含较大质量的物资与具有紧迫时限的货箱，31个硬时限箱必须优先进入可行性判断。由模型可知，单点最大安全载荷受额定质量、爬升能耗与返航余量共同限制，实际批次还可能受体积主导；多点调度则必须在逐站卸载和共享电池再可用之间进行时序协调。

通信约束使中继在位时间成为运输可行性的组成部分。任务分区进一步放大资源不可跨组共享的影响，而这种影响只能沿实际保障关系计算，不能用同波次关系替代。本稿已经给出上述问题的数学表达与检验路径，但最终载荷、调度指标、中继方案、分区配置以及敏感性数值尚需统一复算后补充，不据方法描述推断任何最终最优方案。
